\documentclass[12pt,a4paper]{article}% 文档格式
\usepackage{ctex,hyperref}% 输出汉字
\usepackage{times}% 英文使用Times New Roman

\title{
  \centering % 保证标题居中
  \heiti % 黑体
  \fontsize{22pt}{33pt}\selectfont % 二号字体，1.5倍行距
  改革开放访谈报告 \\[1em] % 主标题
  \hfill % 在下一行之前填充空白
  \fontsize{14pt}{21pt}\selectfont % 副标题字体设置
  \fangsong % 副标题字体样式
  ———改革开放对中国社会和文化的影响 % 副标题
}


\author{
  \centering % 保证作者部分也居中
  \fangsong % 仿宋字体
  \fontsize{14pt}{21pt}\selectfont % 四号字体，1.5倍行距

}

\date{}% 日期（这里避免生成日期）
\usepackage{amsmath,amsfonts,amssymb}% 为公式输入创造条件的宏包
\usepackage{graphicx}% 图片插入宏包
\usepackage{subfigure}% 并排子图
\usepackage{float}% 浮动环境，用于调整图片位置
\usepackage[export]{adjustbox}% 防止过宽的图片
\usepackage{bibentry}
\usepackage{natbib}% 以上2个为参考文献宏包
\usepackage{abstract}% 两栏文档，一栏摘要及关键字宏包
\renewcommand{\abstracttextfont}{\fangsong}% 摘要内容字体为仿宋
\renewcommand{\abstractname}{\textbf{摘\quad 要}}% 更改摘要二字的样式
\usepackage{xcolor}% 字体颜色宏包
\newcommand{\red}[1]{\textcolor[rgb]{1.00,0.00,0.00}{#1}}
\newcommand{\blue}[1]{\textcolor[rgb]{0.00,0.00,1.00}{#1}}
\newcommand{\green}[1]{\textcolor[rgb]{0.00,1.00,0.00}{#1}}
\newcommand{\darkblue}[1]
{\textcolor[rgb]{0.00,0.00,0.50}{#1}}
\newcommand{\darkgreen}[1]
{\textcolor[rgb]{0.00,0.37,0.00}{#1}}
\newcommand{\darkred}[1]{\textcolor[rgb]{0.60,0.00,0.00}{#1}}
\newcommand{\brown}[1]{\textcolor[rgb]{0.50,0.30,0.00}{#1}}
\newcommand{\purple}[1]{\textcolor[rgb]{0.50,0.00,0.50}{#1}}% 为使用方便而编辑的新指令
\usepackage{url}% 超链接
\usepackage{bm}% 加粗部分公式
\usepackage{multirow}
\usepackage{booktabs}
\usepackage{epstopdf}
\usepackage{epsfig}
\usepackage{longtable}% 长表格
\usepackage{supertabular}% 跨页表格
\usepackage{algorithm}
\usepackage{algorithmic}
\usepackage{changepage}% 换页
\usepackage{enumerate}% 短编号
\usepackage{caption}% 设置标题
\captionsetup[figure]{name=\fontsize{10pt}{15pt}\selectfont Figure}% 设置图片编号头
\captionsetup[table]{name=\fontsize{10pt}{15pt}\selectfont Table}% 设置表格编号头

\usepackage{indentfirst}% 中文首行缩进
\usepackage[left=2.50cm,right=2.50cm,top=2.80cm,bottom=2.50cm]{geometry}% 页边距设置
\renewcommand{\baselinestretch}{1.5}% 定义行间距（1.5）
\usepackage{fancyhdr} %设置全文页眉、页脚的格式
\pagestyle{fancy}
\hypersetup{colorlinks=true,linkcolor=black}% 去除引用红框，改变颜色



\begin{document}% 以下为正文内容
	\maketitle% 产生标题，没有它无法显示标题
	\lhead{}% 页眉左边设为空
	\chead{}% 页眉中间设为空
	\rhead{}% 页眉右边设为空
	\lfoot{}% 页脚左边设为空
	\cfoot{\thepage}% 页脚中间显示页码
	\rfoot{}% 页脚右边设为空

  \newpage% 从新的一页继续
	\begin{abstract}
		\fangsong 
    本访谈报告探讨了改革开放对中国社会、经济和文化的深远影响。
    通过与不同年龄段和职业背景的身边的普通人对话，以小见大，从而展示了改革开放前后中国社会的巨大变迁。
    改革开放前，物资短缺、经济封闭，文化生活单调匮乏；
    而随着改革开放的推进，市场经济迅速发展，
    个人和家庭的生活水平显著提升，思想也更加开放，文化多样性增强。
  大家的分享显示了改革开放如何改变了个人的生活轨迹，
    同时也展望了未来中国社会继续发展的可能性。
	\end{abstract}
	
	\begin{adjustwidth}{1.06cm}{1.06cm}
		\fontsize{10.5pt}{15.75pt}\selectfont{\heiti{关键词：}\fangsong{
      改革开放
    、社会变迁
    、市场经济
    、文化多样性
    、生活水平}}
		\bigskip
	\end{adjustwidth}

    %页码格式
	\thispagestyle{empty}

	% % 目录
	% \newpage
	% \begin{center} 
	% 	\tableofcontents
	% \end{center}
	% \setcounter{page}{0}
  %   % 设置当前页面格式 empty表示无页码
	% \thispagestyle{empty}

	% \newpage% 从新的一页继续

	%第一段
\newpage% 从新的一页继续
\section*{访谈主题：改革开放对中国社会和文化的深远影响}

\textbf{我}：回顾改革开放这段历史，这场波澜壮阔的变革不仅推动了中国经济的腾飞，
也深刻影响了社会和文化的各个层面。通过与的分享，我们将探讨改革开放如何改变了中国人的生活，尤其是对于不同时代、不同职业的人群，他们的亲身经历和感受。
首先，我们从改革开放之前的生活开始聊起。那时的社会和经济状况是什么样的呢？请大家回忆一下。

\textbf{爷爷（深吸了一口气，目光有些黯淡）}：改革开放之前，日子真的是很艰难。我还记得那时我们家住在乡下，靠种地为生。那时候是一片片的农田，家里就只有几亩地，种植粮食和蔬菜勉强维持生计。所有的东西都得凭票购买，粮票、油票、布票，每个月领到的配给非常有限，什么东西都要精打细算。冬天的时候，孩子们甚至没有足够的棉衣，只能靠几层破旧的布片御寒。回想起来，那段日子真是令人难忘，也让人倍感珍惜现在的生活。

\textbf{奶奶（点头附和）}：是啊，那时候我们要想买点肉、油，得提前几天打听消息，有时甚至要半夜起床去排队。文化生活也很匮乏，我们家连收音机都没有，唯一的娱乐就是晚上大家聚在一起，聊聊天、讲讲故事，这才算是放松的时刻。可那样的聚会，让我们彼此更加亲近，虽然生活艰苦，但人心却是温暖的。

\textbf{父亲（若有所思）}：那个年代，除了经济上的困难，文化生活也几乎是封闭的。家里的电视直到改革开放后才有的，之前大部分时间都只是听广播和看样板戏。像我小时候，最开心的事就是村里放露天电影，每次都能吸引好几百人来看。每一部电影的放映都像是节日，大家围坐在一起，期待着荧幕上闪现的故事，一起欢笑、一起感动。

\textbf{母亲（微笑着说）}：是啊，那时候的日子就是这么单调，尤其是在农村，日常生活全靠务农。没有娱乐设施，邻里之间也很少有聚会活动，大家忙着糊口，根本没有多余的精力去考虑别的。然而，在这种简单的生活中，我们也收获了许多。大家相互帮助，分享生活中的点滴，让彼此的生活不至于孤单。

\textbf{我}：那么，随着改革开放的到来，你们的生活发生了哪些显著的变化？

\textbf{父亲（满怀激情）}：我还记得那时候政策刚刚放宽，允许大家开设私人企业，我就是在那个时候开始尝试经商的。最早的时候，我和几个朋友一起开了家小店铺，主要做一些简单的批发生意，没想到几年下来，竟然慢慢积累了一些财富。改革开放为我们这些普通人提供了前所未有的机会。我觉得这是中国经济腾飞的起点，正是这段经历让我明白，只要努力，就能改变命运。

\textbf{阿姨（眼中闪烁着光芒）}：对我来说，改革开放对文化领域的影响特别大。我从事文艺工作，开放后，我们可以接触到很多国外的文化作品，不再局限于本土的传统戏剧和宣传片。电影、音乐、文学都变得多元化，给了我们创作者更多的灵感和自由。我记得80年代初，看到第一部来自香港的电影，那种文化冲击真是难以形容。它让我意识到，艺术的表达可以有如此多样的形式。

\textbf{父亲（沉思片刻）}：除了经济上的巨大变化，我觉得改革开放还带来了思想上的解放。以前我们都是听从上级指示，生活一成不变，想要做点创新的事也没有空间。改革开放后，国家鼓励创新和个人发展，这让我们年轻一代看到了新的可能性。我和你母亲也是在那个时候决定进城发展，从农村搬到了城市。虽然开始的时候很难，但我们通过自己的努力，逐渐在城市站稳了脚跟，这样的经历让我深刻感受到，改变始于勇敢的尝试。

\textbf{我}：改革开放不仅影响了经济发展，还改变了大家的生活方式和思维方式。那么，从文化层面上看，大家感受到的变化是什么？

\textbf{奶奶（微笑着说）}：我觉得，文化上的变化是潜移默化的。以前我们接触到的文化非常单一，都是以革命故事和样板戏为主。随着改革开放，电视、电影、音乐的种类丰富起来了，特别是我们可以看到一些来自国外的作品。像我当年看到了来自日本和美国的电视剧，真的让我对外面的世界有了新的认识。这种文化上的冲击让我们意识到，世界原来如此多彩。

\textbf{阿姨（点头赞同）}：不仅仅是文艺，日常生活中的变化也非常显著。比如说，年轻人开始追求时尚，不再拘泥于传统的穿着打扮。我记得你大舅妈那时候花了不少钱给自己买了一条牛仔裤，成了村里的时尚先锋，大家都觉得特别新潮。现在想想，那时候的我们正在慢慢接受外来的新事物，也开始寻找属于自己的生活方式。

\textbf{父亲（激动地说）}：是啊，我觉得改革开放最大的变化还是在思想上。以前我们总觉得生活就是按部就班，没有太多选择。改革开放之后，大家开始意识到，个人的努力和创新能够带来真正的改变。无论是经济上的腾飞，还是文化上的多样性，都是因为人们的思想变得更加开放和包容。每一个小小的进步，都是我们共同努力的结果。

\textbf{我}：听到大家的分享，我感觉改革开放不仅改变了生活条件，也让社会更加多元化。接下来，我想请大家谈谈改革开放对家庭生活的影响。你们的家庭是如何适应这些变化的？

\textbf{爷爷（缓缓地说）}：以前家里都是务农，靠天吃饭，收入微薄。改革开放后，我第一个感受到的就是粮食生产的自由化。我们不再受限于生产队的分配，自己种地，自己决定怎么卖，收入一下子提高了不少。最让我记得深刻的是，我们终于能给孩子们买新衣服了，冬天不用担心他们挨冻。这些小小的变化，对我们来说都是莫大的幸福。

\textbf{父亲（感慨地说）}：搬到城市后，我们生活节奏一下子加快了。以前在农村，生活虽然简单但也安稳，进入城市后，面临着巨大的竞争压力。然而，城市里的工作机会更多，社会福利也逐渐完善，尤其是医疗和教育条件，让我们对未来充满了信心。回头看，虽然艰难，但这一切都是值得的。

\textbf{母亲（充满温柔地说）}：我觉得最重要的变化是孩子们的教育。我们那时候只能读到小学或初中，家里的条件不允许继续深造。而现在，改革开放后的政策给了孩子们更多的机会。你和你表姐都能够接受良好的教育，甚至还有机会出国留学。这是我们那代人无法想象的事情，看到你们有这样的未来，我心里真的很欣慰。

\textbf{我}：确实，教育和家庭条件的提升是改革开放给我们带来的巨大红利。那么对于未来，你们对改革开放继续推动社会发展有哪些期待？

\textbf{爷爷（坚定地说）}：我希望国家能够继续坚持改革开放的方针，特别是在农村经济发展方面做更多的改进。虽然现在生活条件已经大大改善，但我觉得还有很多地方可以进一步发展，尤其是在科技和农业结合方面。科技是未来的方向，我们要好好把握。

\textbf{父亲（充满信心）}：我希望未来的中国能够更加开放，特别是在技术创新和全球合作方面。我相信，中国的未来会更加辉煌，而我们每个人都能从中受益。我们的努力，必将推动社会的进步。

\textbf{阿姨（热情地说）}：我希望未来的文化领域更加多元化，同时能够保持我们自己的文化特色。作为文化工作者，我觉得我们在吸收外来文化的同时，也要注重传承和发扬中华文化的精髓。我们要让世界听到中国的声音。

\textbf{表哥（充满期待）}：作为年轻一代，我希望能够在开放的环境中，继续追求自己的梦想。改革开放给了我们无数的机会，我相信，只要我们努力，未来的世界将会更加广阔。年轻人是未来的希望，我们必须勇于探索。

\textbf{我}：非常感谢大家的分享。通过今天的访谈，我们不仅回顾了改革开放带来的巨大变化，也展望了未来的发展前景。无论是在经济、文化还是家庭生活中，改革开放都发挥了不可估量的作用。让我们继续努力，共同为建设更加美好的未来贡献力量！

	\section*{总结}

  在中国的历史长河中，改革开放是一个具有里程碑意义的时代。从1978年起，这场波澜壮阔的变革彻底改变了中国的经济、社会和文化面貌。在本次访谈中，大家通过分享他们的经历，为我们呈现了改革开放前后鲜明的生活对比，帮助我们更深入地理解这段历史带来的巨大变化。

  访谈伊始，爷爷回忆起改革开放前的生活，那段艰辛的岁月让人感慨不已。凭票购买粮食和日用品的经历，几乎成了当时每一个家庭的共同回忆。在那个年代，农村的经济主要依赖农业，家庭收入微薄，生活条件相当拮据。爷爷讲述了自己在冬天时，孩子们没有足够的棉衣，只能靠几层旧布片来御寒，听着让人心疼。奶奶的回忆也同样唤起了许多人的共鸣，那个时候家里没有娱乐设施，文化生活十分匮乏，大家只能通过聚会和聊天来打发时间。父亲提到，村里的露天电影是孩子们最开心的时刻，每当放映时，大家都聚在一起，期待着银幕上的故事。
  
  然而，改革开放的到来为生活带来了翻天覆地的变化。父亲分享了自己如何抓住这个机会，开始了创业之路。他和朋友们在政策放宽后开了一家小店，逐渐积累财富，这一切都是在改革开放的洪流中实现的。阿姨则提到，作为文艺工作者，改革开放让她接触到丰富的外部文化，电影、音乐的多元化让她和其他创作者获得了新的灵感。开放的政策，不仅仅是经济的解放，更是思想的解放。
  
  谈到文化的变化，爷爷提到，随着改革开放，文化的传播和接受变得更加丰富多彩。人们能够接触到许多国外的作品，这让他们对外部世界有了更深的认识。父亲在讨论中也强调了这种思想上的转变，改革开放之后，大家意识到，个人的努力和创新能够带来真正的改变。
  
  在谈到家庭生活的变化时，爷爷提到，改革开放后，家庭的经济条件有了显著改善。通过自主种植和销售，收入的提高让每个家庭都有了更多的选择和自由。父亲则提到，搬到城市后，生活节奏加快，尽管面临巨大的竞争压力，但城市提供了更多的工作机会和社会福利。母亲则强调，孩子们如今有了更多的教育机会，这与过去的情况截然不同。
  
  在展望未来时，大家表达了对改革开放继续推动社会发展的期待。爷爷希望国家能够在农村经济方面做出更多的改进，父亲则希望中国在技术创新和全球合作上能够更加开放。阿姨提到，作为文化工作者，她希望未来的文化领域能够更加多元化，同时不忘传承中华文化的精髓。表哥则希望能够在这样的开放环境中，继续追求自己的梦想。
  
  通过这次访谈，我们不仅回顾了改革开放所带来的巨大变化，也展望了未来的发展前景。改革开放的成功，体现在物质的丰盈和思想的解放上。每个人的努力汇聚成了推动社会前进的力量。未来，中国在保持传统文化根基的同时，将继续拥抱多元化与全球化的趋势，让每一个人都能在这片土地上找到属于自己的位置。
  

  
\end{document}